\section{Theoretical Setup}\label{sec:setup}

\subsection{Tagged threshold products}

For a positive integer $q$, write $[q]=\{1,\ldots,q\}$.  Fix integers
$k\ge2$ and $N\ge2$, and define
\[
X_{k,N}=[k]\times[N],
\qquad
\tau_t(x)=\mathbf 1\{x\ge t\}
\quad(t\in[N+1],\ x\in[N]).
\]
The tagged threshold-product class and the unrestricted improper output space
are
\[
C_{k,N}=\{c_{\boldsymbol t}:\boldsymbol t\in[N+1]^k\},
\qquad
c_{\boldsymbol t}(i,x)=\tau_{t_i}(x),
\qquad
\mathcal H_{k,N}=\{0,1\}^{X_{k,N}}.
\]
No properness, representation, or computational restriction is imposed.

For a distribution $P$ on $X_{k,N}$, a target
$c_{\boldsymbol t}\in C_{k,N}$, and $h\in\mathcal H_{k,N}$, let
\[
R_P(h,c_{\boldsymbol t})
=\Pr_{Z\sim P}\{h(Z)\ne c_{\boldsymbol t}(Z)\}.
\]
Write $P^{c_{\boldsymbol t}}$ for the realizable labeled law of
$(Z,c_{\boldsymbol t}(Z))$ with $Z\sim P$.  Thus
$S\sim(P^{c_{\boldsymbol t}})^n$ always denotes exactly $n$ independent
labeled examples.

A candidate learner is an arbitrary randomized map
\[
A:(X_{k,N}\times\{0,1\})^n\longrightarrow\mathcal H_{k,N}.
\]
Two ordered size-$n$ labeled datasets are adjacent, denoted $S\simeq S'$,
when they differ in at most one row by replacement.  Every unadorned $\log$
is natural, and
\[
\log_2^*u
=\min\{j\in\mathbb Z_{\ge0}:\log_2^{(j)}u\le1\},
\qquad
m_{n,k}=\max\left\{8,\left\lceil\frac{4n}{k}\right\rceil\right\}.
\]
We use the tower convention
\[
\operatorname{Tow}_2(1)=2,
\qquad
\operatorname{Tow}_2(r+1)=2^{\operatorname{Tow}_2(r)}.
\]

\subsection{Assumptions}

The constants in the following assumptions are fixed before the candidate
parameters and learner are chosen.

\begin{assumption}[Primitive candidate regime]\label{assump:candidate-regime}
For fixed $c_\delta,\varepsilon_0>0$ and $N_0\in\mathbb Z_{\ge2}$, the
candidate parameters satisfy
\[
k\ge2,
\qquad N\ge N_0,
\qquad n\in\mathbb Z_{\ge1},
\qquad 0<\varepsilon\le\varepsilon_0,
\]
and, with $M=m_{n,k}$,
\[
0<\delta\le
\min\left\{
\frac{1}{n\log(n+1)},
\frac{c_\delta}{M^2\log(M+1)}
\right\}.
\]
\end{assumption}

\begin{assumption}[Central approximate differential privacy]\label{assump:central-dp}
For every $S\simeq S'$ in $(X_{k,N}\times\{0,1\})^n$ and every event
$E\subseteq\mathcal H_{k,N}$,
\[
\Pr\{A(S)\in E\}
\le e^\varepsilon\Pr\{A(S')\in E\}+\delta,
\]
where the probabilities are over only the internal randomness of $A$.
\end{assumption}

\begin{assumption}[Distribution-free realizable PAC guarantee]\label{assump:distribution-free-realizable-pac}
For fixed $\alpha_0,\beta_0\in(0,1/2)$, every
$\boldsymbol t\in[N+1]^k$ and every distribution $P$ on $X_{k,N}$ satisfy
\[
\Pr_{\substack{S\sim(P^{c_{\boldsymbol t}})^n\\A}}
\left\{R_P(A(S),c_{\boldsymbol t})\le\alpha_0\right\}
\ge1-\beta_0.
\]
The probability includes the exact i.i.d. sample and the learner's internal
randomness.
\end{assumption}

\subsection{Objective}

We seek a deterministic, pointwise-in-$n$ lower-bound implication.  Its
probabilistic premise is the fixed-size PAC guarantee above; its conclusion
does not average over candidate sample sizes, targets, or distributions.  The
loss throughout is population zero-one risk.
